% =============================================================================
% RCS --- Formal Definitions
% Recursive Controlled Systems: Core mathematical framework
% =============================================================================
\documentclass[11pt,a4paper]{article}

% --- Encoding and fonts ---
\usepackage[utf8]{inputenc}
\usepackage[T1]{fontenc}
\usepackage{lmodern}

% --- Language ---
\usepackage[english]{babel}

% --- Page geometry ---
\usepackage[margin=2.5cm]{geometry}

% --- Colors (must load before hyperref for color mixing) ---
\usepackage{xcolor}

% --- Hyperlinks ---
\usepackage[colorlinks=true,linkcolor=blue!70!black,citecolor=blue!70!black,urlcolor=blue!70!black]{hyperref}

% --- Tables ---
\usepackage{booktabs}
\usepackage{array}

% --- Paragraph spacing ---
\usepackage{parskip}

% --- Micro-typography ---
\usepackage{microtype}

% --- Shared RCS preamble ---
% \SHAREDLATEX points at the shared LaTeX assets dir. Each paper in
% papers/pN-.../ resolves it relative to its own location; preamble.tex
% uses the macro to reference its own includes (parameters.tex, etc.).
\def\SHAREDLATEX{../../latex}
% =============================================================================
% RCS — Shared LaTeX Preamble
% Used by all papers in the Recursive Controlled Systems series
% =============================================================================

% --- Math ---
\usepackage{amsmath,amssymb,amsthm}
\usepackage{mathtools}
\usepackage{bm}              % Bold math symbols

% --- Theorem environments ---
\theoremstyle{plain}
\newtheorem{theorem}{Theorem}[section]
\newtheorem{lemma}[theorem]{Lemma}
\newtheorem{proposition}[theorem]{Proposition}
\newtheorem{corollary}[theorem]{Corollary}

\theoremstyle{definition}
\newtheorem{definition}[theorem]{Definition}
\newtheorem{example}[theorem]{Example}
\newtheorem{remark}[theorem]{Remark}
\newtheorem{assumption}[theorem]{Assumption}
\newtheorem{law}{Law}

% --- RCS notation macros ---

% Sets and spaces
\newcommand{\X}{\mathcal{X}}           % State space
\newcommand{\Y}{\mathcal{Y}}           % Observation space
\newcommand{\U}{\mathcal{U}}           % Control input space
\newcommand{\M}{\mathcal{M}}           % Mutation operator set
\newcommand{\setR}{\mathbb{R}}         % Real numbers
\newcommand{\setN}{\mathbb{N}}         % Natural numbers
\newcommand{\setZ}{\mathbb{Z}}         % Integers

% RCS 7-tuple
\newcommand{\rcs}{\Sigma}              % A Recursive Controlled System
\newcommand{\plant}{\rcs}              % Alias: plant
\newcommand{\dyn}{f}                   % Dynamics function
\newcommand{\obs}{h}                   % Observation map
\newcommand{\shield}{S}                % Safety shield
\newcommand{\ctrl}{\Pi}                % Controller

% Levels
\newcommand{\Li}[1]{L_{#1}}           % Level i: \Li{0}, \Li{1}, etc.

% Stability budget
\newcommand{\stab}{\lambda}            % Stability margin
\newcommand{\decay}{\gamma}            % Nominal decay rate
\newcommand{\adapt}{L_\theta \rho}     % Adaptation cost
\newcommand{\design}{L_d \eta}         % Design evolution cost
\newcommand{\delay}{\beta \bar{\tau}}  % Delay cost
\newcommand{\switch}{\frac{\ln \nu}{\tau_a}} % Switching cost

% Homeostatic drive
\newcommand{\drive}[1]{D_{#1}}         % Drive at level i: \drive{i}
\newcommand{\setpoint}[1]{x^*_{#1}}   % Setpoint at level i
\newcommand{\homostate}{\xi}           % Homeostatic state vector

% Lyapunov
\newcommand{\lyap}[1]{V_{#1}}         % Lyapunov function at level i
\newcommand{\dlyap}{\dot{V}}          % Time derivative of V

% EGRI
\newcommand{\egri}{\Pi_{\text{EGRI}}}  % EGRI meta-controller
\newcommand{\budget}{B}                % Budget
\newcommand{\score}{J}                 % Evaluator/score function
\newcommand{\trial}{\mathcal{T}}       % Trial record

% Agents and information state
\newcommand{\info}{I}                  % Information state (Eslami)
\newcommand{\mem}{m}                   % Memory state
\newcommand{\toolout}{z}               % Tool output
\newcommand{\interact}{r}              % Interaction signal
\newcommand{\param}{\theta}            % Adaptable parameters
\newcommand{\goal}{\zeta}              % Goal descriptor
\newcommand{\mode}{\sigma}             % Mode/strategy index
\newcommand{\arch}{c}                  % Architecture/workflow config

% Eslami delays
\newcommand{\taubar}{\bar{\tau}}       % Supremal total delay
\newcommand{\tauu}{\tau_u}             % Control computation delay
\newcommand{\tauth}{\tau_\theta}       % Adaptation update delay
\newcommand{\tauz}{\tau_z}             % Tool/retrieval delay
\newcommand{\taus}{\tau_\sigma}        % Switching decision delay
\newcommand{\tauc}{\tau_c}             % Reconfiguration delay
\newcommand{\taua}[1][]{\tau_{a\ifx\\#1\\\else,#1\fi}} % Average dwell time: \taua = tau_a, \taua[i] = tau_{a,i}

% Norms
\newcommand{\norm}[1]{\left\| #1 \right\|}
\newcommand{\abs}[1]{\left| #1 \right|}

% Operators
\DeclareMathOperator*{\argmin}{arg\,min}
\DeclareMathOperator*{\argmax}{arg\,max}


% --- Bibliography ---
\usepackage[numbers,sort&compress]{natbib}

% =============================================================================
\title{Recursive Controlled Systems:\\
Formal Definitions and Core Theorems}
\author{Carlos D.\ Escobar-Valbuena}
\date{April 2026 --- Working Draft}

\begin{document}
\maketitle

\begin{abstract}
We introduce the \emph{Recursive Controlled System} (RCS), a control-theoretic
\emph{type signature} $(\X, \Y, \U, \dyn, \obs, \shield, \ctrl)$ whose
controller $\ctrl$ is itself an RCS at the next hierarchical level.  The
formalization is agnostic with respect to the controller class: classical
regulators, latent world models, JEPA-family predictors, and LLM-based
policies each enter as instantiations that discharge the same assumptions on
their own state space.  We define the homeostatic drive as a common notation
covering Lyapunov, reward, and free-energy interpretations in their
respective regimes.  Extending the stability budget of Eslami \&
Yu~\cite{eslami2026control} to finite-depth hierarchies, we prove by
induction on depth~$N$ that if the per-level stability budget is positive
and a time-scale-separation bound holds, the composite state decays
exponentially at rate $\omega_N = \min_i \stab_i$.  An instantiation
catalogue illustrates how LQR, Dreamer, V-JEPA~2, and LLM controllers plug
into the type signature.
\end{abstract}

\tableofcontents
\newpage

% =============================================================================
\section{Introduction}
\label{sec:introduction}
% =============================================================================

Autonomous AI agents increasingly operate as controllers of complex
systems, from energy microgrids to software infrastructure.  Yet the
agent itself~--- its internal state, cognitive resources, and economic
constraints~--- is a dynamical system that must also be regulated.
This creates a recursive structure: one controller regulates a plant, a
meta-controller regulates the first controller, and a governance layer
constrains the meta-controller.

\paragraph{Framework, not instantiation.}
We formalize this recursive structure \emph{agnostically with respect to
the controller implementation}.  The Recursive Controlled System (RCS)
is a 7-tuple
$(\X, \Y, \U, \dyn, \obs, \shield, \ctrl)$~--- a \emph{type signature}
over which stability theorems operate~--- not a commitment to any
particular plant or policy class.  Concrete controllers (classical LQR
regulators~\cite{khalil2002nonlinear}, Dreamer-style latent world
models~\cite{hafner2025dreamerv3}, JEPA-family
predictors~\cite{assran2025vjepa2}, and LLM-based
policies~\cite{bhargava2023magic,soatto2023taming}) are
\emph{instantiations} of the type signature, each entering the
framework by discharging the per-level assumptions of
Section~\ref{sec:stability} in its own regime.  This
architecture-agnostic stance is what the main theorem requires: what
matters is not what $\ctrl$ is, but whether the per-level stability
budget is positive.  Section~\ref{sec:instantiations} catalogues four
representative instantiations with their stability witnesses.

\paragraph{Open problem this paper addresses.}
The closest published articulation of the same ambition is LeCun's
position paper~\cite{lecun2022path}, which sketches a six-module
autonomous agent architecture with an intrinsic cost module whose
stability is \emph{assumed} rather than derived, and an informal
two-level hierarchical JEPA (H-JEPA) whose composite stability is not
stated as a theorem.  Our finite-depth induction
(Theorem~\ref{thm:recursive}) closes this gap: given per-level
assumptions that modern world models~\cite{hafner2025dreamerv3,
assran2025vjepa2}, classical regulators~\cite{khalil2002nonlinear}, and
LLM controller architectures~\cite{bhargava2023magic,soatto2023taming,
wang2026agentspec} are each known to satisfy on their respective state
spaces, we obtain composite exponential stability for any hierarchy of
finite depth.

\paragraph{Contributions.}
\begin{enumerate}
  \item The RCS definition as a recursive \emph{type signature}, the
        same 7-tuple re-applied at every level
        (Section~\ref{sec:rcs}).
  \item The homeostatic drive as a common notation covering three
        regime-restricted interpretations~--- Lyapunov, reward,
        free-energy bound~--- with each regime stated explicitly
        (Section~\ref{sec:drive}).
  \item A recursive stability budget extending Eslami \&
        Yu~\cite{eslami2026control} Theorem~2 to finite-depth
        hierarchies, proven by induction on depth
        (Section~\ref{sec:stability}).
  \item An instantiation catalogue showing how LQR, Dreamer, V-JEPA~2,
        and LLM-as-controller each discharge the assumptions
        (Section~\ref{sec:instantiations}).
  \item The EGRI meta-controller as an instance of $\rcs_2$, with two
        of its Five Laws admitting derived formal content and the
        remaining three recast as design principles
        (Section~\ref{sec:egri}).
\end{enumerate}


% =============================================================================
\section{Preliminaries}
% =============================================================================

We recall standard definitions from nonlinear systems theory, following
Khalil~\cite{khalil2002nonlinear} and Liberzon~\cite{liberzon2003switching}.

\begin{definition}[Discrete-time dynamical system]
A \emph{discrete-time dynamical system} is a triple $(\X, \U, \dyn)$ where
$\X \subseteq \setR^n$ is the state space, $\U \subseteq \setR^m$ is the
input space, and $\dyn : \X \times \U \to \X$ is the state transition map:
\begin{equation}
  x_{k+1} = \dyn(x_k, u_k).
\end{equation}
\end{definition}

\begin{definition}[Lyapunov stability]
An equilibrium $x^*$ of $x_{k+1} = \dyn(x_k, 0)$ is \emph{exponentially
stable} if there exist constants $K > 0$ and $0 < \omega < 1$ such that
\begin{equation}
  \norm{x_k - x^*} \leq K \, \omega^k \, \norm{x_0 - x^*}
  \quad \forall\, k \geq 0.
\end{equation}
\end{definition}

\begin{definition}[Average dwell time~{\cite{hespanha1999stability}}]
\label{def:dwell-time}
A switching signal $\mode(t)$ has \emph{average dwell time} $\taua$ if
\begin{equation}
  N_\mode(t, t+T) \leq N_0 + \frac{T}{\taua}
  \quad \forall\, t \geq 0,\; T \geq 0,
\end{equation}
where $N_\mode(t, t+T)$ is the number of switches in $[t, t+T]$.
\end{definition}

\begin{remark}[Standing-max convention for $N_0$ across levels]
\label{rem:N0-standing-max}
In the recursive setting (Theorem~\ref{thm:recursive}), we write
$N_0 := \max_i N_{0,i}$, where $N_{0,i}$ is the chatter bound of
level~$i$'s switching signal; the hypothesis (H5) is to be read under
this convention.  Writing the shared constant $N_0$ is then
mathematically equivalent to taking the (finite) per-level supremum
and keeps the recursion statement uniform.
\end{remark}


% =============================================================================
\section{Controlled System}
\label{sec:controlled-system}
% =============================================================================

\begin{definition}[Controlled System]
\label{def:cs}
A \emph{Controlled System} is a 7-tuple
\begin{equation}
  \rcs = (\X, \Y, \U, \dyn, \obs, \shield, \ctrl)
\end{equation}
where:
\begin{itemize}
  \item $\X$ is the \emph{state space} --- the set of possible internal states;
  \item $\Y$ is the \emph{observation space} --- the set of possible measurements;
  \item $\U$ is the \emph{control input space} --- the set of possible actions;
  \item $\dyn : \X \times \U \to \X$ is the \emph{dynamics} --- the state
        transition function;
  \item $\obs : \X \to \Y$ is the \emph{observation map} --- how state becomes
        measurement;
  \item $\shield : \U \times \X \to \U$ is the \emph{safety shield} --- a map
        that filters any proposed control input to a safe one:
        \begin{equation}
          u_{\text{safe}} = \shield(\tilde{u}, x)
          = \argmin_{u \in \U_{\text{safe}}(x)} \norm{u - \tilde{u}}^2;
        \end{equation}
  \item $\ctrl : \Y^* \to \U$ is the \emph{controller} --- a map from
        observation histories to control inputs.
\end{itemize}

Here $\Y^* = \bigcup_{k=0}^{\infty} \Y^k$ denotes the set of finite
observation sequences.  In practice, $\ctrl$ maintains internal state --- this
is formalized in the recursive definition below.
\end{definition}

\begin{remark}[Shield as CBF-QP]
When $\U_{\text{safe}}(x) = \{u \in \U : \nabla_x h(x) \cdot \dyn(x, u)
+ \alpha(h(x)) \geq 0\}$ for a control barrier function $h : \X \to \setR$
with class-$\mathcal{K}$ function $\alpha$, the shield reduces to the
standard CBF-QP formulation~\cite{eslami2026control}.
\end{remark}


% =============================================================================
\section{Recursive Controlled System}
\label{sec:rcs}
% =============================================================================

\begin{definition}[Recursive Controlled System]
\label{def:rcs}
A \emph{Recursive Controlled System} (RCS) is a Controlled System
$\rcs = (\X, \Y, \U, \dyn, \obs, \shield, \ctrl)$ where the controller
$\ctrl$ is itself a Controlled System at the next hierarchical level:
\begin{equation}
  \ctrl = \rcs' = (\X', \Y', \U', \dyn', \obs', \shield', \ctrl').
\end{equation}

This produces a hierarchy of levels $\Li{0}, \Li{1}, \Li{2}, \ldots$ where:
\begin{enumerate}
  \item Level $\Li{i}$ is a Controlled System whose controller is $\Li{i+1}$.
  \item The state $\X_{i+1}$ includes the parameters and configuration of
        $\Li{i}$'s controller.
  \item The observation $\Y_{i+1}$ is derived from the performance metrics of
        $\Li{i}$.
  \item The control output $\U_{i+1}$ modifies the behavior of $\Li{i}$.
\end{enumerate}

We write $\rcs_i = (\X_i, \Y_i, \U_i, \dyn_i, \obs_i, \shield_i, \ctrl_i)$
for the Controlled System at level $\Li{i}$.
\end{definition}

\begin{remark}[Structural self-similarity]
\label{rem:structural-self-similarity}
The hierarchy uses the same 7-tuple signature at every level: at
level~$\Li{i}$, the controller slot $\ctrl_i$ is itself a Controlled
System at level~$\Li{i+1}$.  In practice, the hierarchy has finite
depth $N \leq 4$ (plant, adapter, meta-controller, governance); the
recursion is notational convenience, and
Theorem~\ref{thm:recursive} proves composite stability by induction on
this finite depth, without invoking a categorical fixed point.
\end{remark}


% =============================================================================
\section{Homeostatic Drive as Lyapunov Function}
\label{sec:drive}
% =============================================================================

\begin{definition}[Homeostatic drive]
\label{def:drive}
At each level $\Li{i}$, the \emph{homeostatic drive} is:
\begin{equation}
  \drive{i}(x) = \norm{x - \setpoint{i}}^2, \quad x \in \X_i,
\end{equation}
where $\setpoint{i} \in \X_i$ is the \emph{setpoint} (desired equilibrium) at
level $i$.
\end{definition}

\begin{proposition}[Three regime-restricted interpretations of $\drive{i}$]
\label{prop:triple}
The drive $\drive{i}$ admits three interpretations, each holding in its
own regime:
\begin{enumerate}
  \item \textbf{Lyapunov function (quadratic-$V$ regime).}  Whenever
        $\drive{i}$ satisfies the local sandwich of
        Assumption~\ref{ass:per-level}, strict decrease
        $\drive{i}(x_{k+1}) < \drive{i}(x_k)$ for all $x_k \neq
        \setpoint{i}$ certifies asymptotic stability of $\setpoint{i}$
        in the standard Lyapunov
        sense~\cite{khalil2002nonlinear}.
  \item \textbf{Reward signal (bounded physiological drives).}  For
        pre-specified homeostatic setpoints, the drive-reduction reward
        $r_k = \drive{i}(x_k) - \drive{i}(x_{k+1})$ recovers the
        homeostatic reinforcement learning
        framework~\cite{keramati2014homeostatic}.  Extending this
        identification to arbitrary RL reward shapings requires
        additional hypotheses and is not claimed here.
  \item \textbf{Free-energy bound (linear-Gaussian regime).}  Under a
        linear-Gaussian generative model with prior preferences centred
        at $\setpoint{i}$, $\drive{i}$ upper-bounds the variational
        free energy, recovering active
        inference~\cite{friston2009reinforcement,baltieri2019pid}.
\end{enumerate}
The three regimes coincide exactly when both of the following hold:
\begin{itemize}
  \item[\textbf{(A)}] \emph{Linear-quadratic structure.}  The dynamics
        are linear-in-state, $\dot{x} = Ax + Bu$, and the cost is
        quadratic, $\ell(x,u) = x^\top Q x + u^\top R u$ with
        $Q \succeq 0$ and $R \succ 0$.
  \item[\textbf{(B)}] \emph{Fixed setpoint.}  The setpoint
        $\setpoint{i}$ is pre-specified and held constant, not adapted
        online.
\end{itemize}
Under (A)+(B), the three readings relate as follows.
\emph{Lyapunov $\Leftrightarrow$ reward}: the quadratic cost admits
a value function $V(x) = x^\top P x$ where $P$ solves the algebraic
Riccati equation~\cite{khalil2002nonlinear}; the drive-reduction
reward $r_k = V(x_k) - V(x_{k+1})$ is the value decrement, so $V$ is
simultaneously the Lyapunov certificate and the (negated) cost-to-go.
\emph{Lyapunov $\Leftrightarrow$ free-energy}: under (A)+(B) the
Gaussian prior on trajectories makes the variational free-energy
bound tight; it equals $V$ up to a log-partition constant, so strict
decrease of $\drive{i}$ coincides with decrease of free
energy~\cite{baltieri2019pid}.  \emph{Reward $\Leftrightarrow$
free-energy} follows by composition.  Outside (A)+(B) the three
readings diverge: a single notation unifies them, a single theorem
does not.
\end{proposition}

\begin{remark}[Bounded objectives prevent unbounded resource acquisition]
Unlike unbounded reward maximization, the homeostatic drive saturates at
$\drive{i}(\setpoint{i}) = 0$.  Once equilibrium is reached, the
controller has no gradient to acquire additional resources, which
prevents \emph{unbounded resource acquisition} as a failure
mode~\cite{mineault2024neuroai}.  This is narrower than ``safety'' in
the specification-gaming or Goodharting sense: bounded objectives still
admit reward hacking, distributional shift, and side-effect failures,
all of which require separate analysis outside the scope of this paper.
\end{remark}


% =============================================================================
\section{Recursive Stability Budget}
\label{sec:stability}
% =============================================================================

We extend the stability budget of Eslami \& Yu~\cite{eslami2026control}
(Theorem~2) to the recursive setting.

\begin{assumption}[Per-level Lyapunov functions]
\label{ass:per-level}
For each level $\Li{i}$, there exists a continuously differentiable Lyapunov
function $\lyap{i} : \X_i \to \setR_{\geq 0}$ and constants
$\underline{\alpha}_i, \bar{\alpha}_i > 0$ such that:
\begin{equation}
  \underline{\alpha}_i \norm{x}^2
  \leq \lyap{i}(x)
  \leq \bar{\alpha}_i \norm{x}^2
  \quad \forall\, x \in \X_i.
\end{equation}
\end{assumption}

\begin{assumption}[Frozen decay]
\label{ass:decay}
Along any flow interval in which all higher-level variables are held constant,
the dynamics at level $\Li{i}$ satisfy:
\begin{equation}
  \lyap{i}(x_{k+1}) - \lyap{i}(x_k) \leq -\decay_i \, \lyap{i}(x_k)
\end{equation}
for some $\decay_i > 0$.
\end{assumption}

\begin{assumption}[Adaptation sensitivity]
\label{ass:adapt}
There exist constants $L_{\param,i}, \rho_i \geq 0$ such that the perturbation
from level $\Li{i+1}$ adjusting parameters $\param_i$ satisfies:
\begin{equation}
  \abs{\frac{\partial \lyap{i}}{\partial \param} \cdot \dot{\param}_i}
  \leq L_{\param,i} \, \rho_i \, \lyap{i}(x_k),
  \quad
  \norm{\dot{\param}_i} \leq \rho_i.
\end{equation}
\end{assumption}

\begin{assumption}[Design evolution sensitivity]
\label{ass:design}
There exist constants $L_{d,i}, \eta_i \geq 0$ such that the perturbation from
level $\Li{i+1}$ changing architecture/design variables $d_i$ satisfies:
\begin{equation}
  \abs{\frac{\partial \lyap{i}}{\partial d} \cdot \dot{d}_i}
  \leq L_{d,i} \, \eta_i \, \lyap{i}(x_k),
  \quad
  \norm{\dot{d}_i} \leq \eta_i.
\end{equation}
\end{assumption}

\begin{assumption}[Delay penalty]
\label{ass:delay}
The total effective delay $\tau_{\text{tot},i} \leq \taubar_i$ introduces a
Lyapunov decrease penalty bounded by:
\begin{equation}
  \Delta_\tau \lyap{i} \leq \beta_i \, \taubar_i \, \lyap{i}(x_k)
\end{equation}
for some $\beta_i \geq 0$.
\end{assumption}

\begin{assumption}[Jump comparability]
\label{ass:jump}
At each mode-switching instant at level $\Li{i}$, there exists $\nu_i \geq 1$
such that:
\begin{equation}
  \lyap{i}(x_k^+) \leq \nu_i \, \lyap{i}(x_k^-).
\end{equation}
\end{assumption}

\begin{definition}[Stability budget at level $\Li{i}$]
\label{def:budget}
The \emph{stability budget} at level $\Li{i}$ is:
\begin{equation}
  \boxed{
  \stab_i
  = \decay_i
  - L_{\param,i} \, \rho_i
  - L_{d,i} \, \eta_i
  - \beta_i \, \taubar_i
  - \frac{\ln \nu_i}{\taua[i]}
  }
  \label{eq:budget}
\end{equation}
where:
\begin{center}
\begin{tabular}{@{}cll@{}}
  \toprule
  \textbf{Term} & \textbf{Name} & \textbf{Source} \\
  \midrule
  $\decay_i$ & Nominal decay rate & Plant dynamics at level $i$ \\
  $L_{\param,i} \rho_i$ & Adaptation cost & Level $i{+}1$ tuning level $i$ parameters \\
  $L_{d,i} \eta_i$ & Design cost & Level $i{+}1$ changing level $i$ architecture \\
  $\beta_i \taubar_i$ & Delay cost & Inference, tool, communication latency \\
  $\frac{\ln \nu_i}{\taua[i]}$ & Switching cost & Mode transitions at level $i$ \\
  \bottomrule
\end{tabular}
\end{center}
\end{definition}

\begin{remark}[Context collapse as a design-evolution violation]
\label{rem:context-collapse}
Empirical evidence for Assumption~\ref{ass:design} comes from the
context-engineering literature.  \citet{zhang2025ace} document a
\emph{context collapse} event in which a monolithic rewrite of an
accumulated context (their Fig.~2, baseline trajectory) causes a
single-step drop from $18{,}282$ to $122$ tokens, with accuracy
falling from $66.7\%$ to $57.1\%$~--- below the no-context baseline
of $63.7\%$.  This is a design-evolution perturbation with
$\eta_1$-ratio $\approx 150\times$ in one step, far exceeding the
admissible bound
\begin{equation}
  \eta_1 <
  \frac{\decay_1 - L_{\param,1}\rho_1 - \beta_1\taubar_1
        - \ln\nu_1/\taua[1]}{L_{d,1}}
\end{equation}
implied by~\eqref{eq:budget}.  Append-only update disciplines
(Reflexion~\cite{shinn2023reflexion}, Voyager~\cite{wang2023voyager},
Self-Refine~\cite{madaan2023selfrefine}, ACE~\cite{zhang2025ace})
automatically satisfy this bound by capping $\eta_1$ at the per-sample
admission rate; monolithic-rewrite systems do not.  This gives a
formal explanation for why the empirical trend across this lineage
converges on incremental, structured updates.
\end{remark}


\paragraph{Standing hypotheses for Theorem~\ref{thm:recursive}.}
For a finite RCS hierarchy $\rcs_0, \rcs_1, \ldots, \rcs_N$ with
$\ctrl_i = \rcs_{i+1}$ for $i < N$, take the following as standing
hypotheses throughout Section~\ref{sec:stability}:
\begin{description}
  \item[(H1) Per-level Lyapunov.]  Each level admits $\lyap{i}$ per
        Assumption~\ref{ass:per-level}.
  \item[(H2) Frozen decay.]  With higher levels held constant,
        Assumption~\ref{ass:decay} gives $\decay_i > 0$ for each $i$.
  \item[(H3) Adaptation coupling.]  Assumption~\ref{ass:adapt} provides
        $L_{\param,i}$ and $\rho_i$ for each $i < N$; at $i = N$ the
        term is vacuous.
  \item[(H4) Design coupling.]  Assumption~\ref{ass:design} provides
        $L_{d,i}$ and $\eta_i$ for each $i < N$; at $i = N$ the term is
        vacuous.
  \item[(H5) Jumps and delay.]  Assumptions~\ref{ass:delay}
        and~\ref{ass:jump} provide $\beta_i$, $\taubar_i$, $\nu_i$,
        $\taua[i]$ for each level.
  \item[(H6) Time-scale separation.]  There exists a constant $c > 1$
        such that for every $i < N$,
        \begin{equation}
          \taua[i+1] \geq c \cdot \taua[i].
          \label{eq:dwell-separation}
        \end{equation}
  \item[(H7) Positive stability budget.]  For every
        $i \in \{0, 1, \ldots, N\}$, $\stab_i > 0$.
\end{description}
The separation constant $c$ in~\eqref{eq:dwell-separation} may be
arbitrarily close to~$1$; the proof relies on strict separation ($c > 1$)
but no specific ratio.

\begin{theorem}[Recursive stability, inductive form]
\label{thm:recursive}
Under (H1)--(H7), any finite RCS hierarchy of depth $N \geq 0$ admits a
composite Lyapunov function
$\lyap{\leq N} : \prod_{i=0}^{N} \X_i \to \setR_{\geq 0}$
and constants $K_N > 0$, $\omega_N > 0$ such that, for every trajectory
$\chi(t) = (x_0(t), x_1(t), \ldots, x_N(t))$,
\begin{equation}
  \lyap{\leq N}(\chi(t))
  \leq K_N \, e^{-\omega_N t} \, \lyap{\leq N}(\chi(0)),
  \qquad \omega_N = \min_{0 \leq i \leq N} \stab_i > 0.
\end{equation}
In particular, the composite state decays exponentially at rate
$\omega_N$.
\end{theorem}

\begin{proof}
Let $P(N)$ denote the conclusion of Theorem~\ref{thm:recursive} for
hierarchies of depth $\leq N$.  We prove $P(N)$ holds for all
$N \geq 0$ by induction on~$N$.

\emph{Base case ($N = 0$).}  A depth-$0$ hierarchy is a single
Controlled System $\rcs_0$ with no level above it, so (H3) and (H4)
are vacuous at $i = 0$.  By (H1) at level~$0$, $\lyap{0}$ sandwiches
$\norm{x_0 - \setpoint{0}}^2$.  By (H2) at level~$0$, along flow
intervals
$\lyap{0}(x^{k+1}) - \lyap{0}(x^k) \leq -\decay_0 \lyap{0}(x^k)$.
By (H5) at level~$0$ (delay and jump contributions),
$\Delta \lyap{0} \leq -(\stab_0 + \ln\nu_0/\taua[0]) \lyap{0}$ between
jumps and $\lyap{0}(x^+) \leq \nu_0 \lyap{0}(x^-)$ at jumps.  Combining
flow decay with jump growth via the Hespanha--Morse average
dwell-time formula~\cite{hespanha1999stability} gives
\begin{equation}
  \lyap{0}(t)
  \leq \nu_0^{N_0} \, e^{-\stab_0 t} \, \lyap{0}(0),
  \label{eq:base-decay}
\end{equation}
where $N_0$ is the dwell-time chatter bound.  Setting
$\lyap{\leq 0} := \lyap{0}$, $\omega_0 := \stab_0$,
$K_0 := \nu_0^{N_0}$ and invoking (H7), we conclude $P(0)$.

\emph{Inductive step.}  Assume $P(n)$ and prove $P(n+1)$.  Consider a
hierarchy of depth $n + 1$: $\rcs_0, \ldots, \rcs_n, \rcs_{n+1}$.

\smallskip
\noindent\textbf{(i)~Sub-hierarchy bound.}  Apply the inductive
hypothesis to the sub-hierarchy $\rcs_0, \ldots, \rcs_n$, treating
$\rcs_{n+1}$ as a quasi-static source of adaptation and design
perturbations entering level~$n$ through (H3) and (H4).  This yields
$\lyap{\leq n}$ with
$\lyap{\leq n}(t) \leq K_n e^{-\omega_n t} \lyap{\leq n}(0)$ and
$\omega_n = \min_{i \leq n} \stab_i$.

\smallskip
\noindent\textbf{(ii)~Top-level bound.}  By (H1), (H2), and (H5) at
level~$n+1$ (with (H3) and (H4) vacuous at the top level since no
$\rcs_{n+2}$ exists), a dwell-time argument identical to the base case
yields
\begin{equation}
  \lyap{n+1}(t)
  \leq \nu_{n+1}^{N_0}\, e^{-\stab_{n+1} t}\, \lyap{n+1}(0).
  \label{eq:top-decay}
\end{equation}

\smallskip
\noindent\textbf{(iii)~Coupling.}  The coupling of $\rcs_{n+1}$ into
the sub-hierarchy is, by construction, the perturbation $\rcs_{n+1}$
delivers to level~$n$ as adaptation (rate $\rho_n$) and design
(rate~$\eta_n$) evolution.  These contributions are already absorbed
into $\stab_n$ via (H3) and (H4), and therefore into $\omega_n$; no new
inter-level coupling is introduced at step~(iii).

\smallskip
\noindent\textbf{(iv)~Time-scale separation.}  By (H6),
$\taua[n+1] \geq c \cdot \taua[n]$ with $c > 1$ places level~$n+1$
strictly slower than the sub-hierarchy top level.  This satisfies the
time-scale hypothesis of Tikhonov's singular-perturbation theorem
(Khalil~\cite{khalil2002nonlinear} Thm.~11.4 / Liberzon
\cite{liberzon2003switching} Thm.~3.1): level~$n+1$ is the slow
subsystem, the depth-$n$ sub-hierarchy is the fast one, and holding
$x_{n+1}$ quasi-static is valid up to an $O(c^{-1})$ error that (H6)
makes arbitrarily small.

\smallskip
\noindent\textbf{(v)~Composite Lyapunov construction.}  Choose
$\alpha_{n+1} > 0$ per the standard Tikhonov construction~\cite{khalil2002nonlinear}
so that cross-term contributions of~$\alpha_{n+1}\lyap{n+1}$ into the
sub-hierarchy are dominated by $\omega_n\lyap{\leq n}$; define
\begin{equation}
  \lyap{\leq n+1}(\chi)
  := \lyap{\leq n}(x_0, \ldots, x_n)
  + \alpha_{n+1} \, \lyap{n+1}(x_{n+1}).
  \label{eq:composite-construction}
\end{equation}
Differentiating along trajectories and applying~\eqref{eq:top-decay}
together with the sub-hierarchy bound from step~(i) gives
\begin{equation}
  \Delta \lyap{\leq n+1}
  \leq -\min(\omega_n, \stab_{n+1}) \cdot \lyap{\leq n+1}
  = -\omega_{n+1} \cdot \lyap{\leq n+1},
  \label{eq:composite-decay}
\end{equation}
with $\omega_{n+1} = \min(\omega_n, \stab_{n+1}) = \min_{i \leq n+1}
\stab_i$.  By (H7), $\omega_{n+1} > 0$.  Taking
$K_{n+1} := K_n \cdot \nu_{n+1}^{N_0}$ completes the inductive step.

By induction, $P(N)$ holds for every $N \geq 0$.
\end{proof}

\begin{remark}[Orphan check: each standing hypothesis is invoked]
\label{rem:orphan-check}
Every standing hypothesis (H1)--(H7) is used at least once in the
proof.  (H1) is used in the base case and step~(ii).  (H2) is used in
the base case and step~(ii).  (H3) and (H4) are used in step~(iii) to
place inter-level adaptation and design perturbations inside $\stab_n$,
and through the inductive hypothesis of step~(i).  (H5) supplies
$\nu_i$, $\taubar_i$, and $\taua[i]$ to the dwell-time bounds
in~\eqref{eq:base-decay} and~\eqref{eq:top-decay}.  (H6) is used in
step~(iv) to satisfy Tikhonov's time-scale hypothesis.  (H7) is used
in the base case to give $\omega_0 > 0$ and in step~(v) to conclude
$\omega_{n+1} > 0$.  No orphan hypotheses remain.
\end{remark}


% =============================================================================
\section{Observer--Controller Lens}
\label{sec:self-similar}
% =============================================================================

The recursion of Definition~\ref{def:rcs} supplies the same 7-tuple
signature at every level.  At each level, the interface between
observation and control admits a \emph{lens} structure, which we record
here for use in Section~\ref{sec:observer}.

\begin{definition}[Observer-controller lens]
\label{def:lens}
The interface between a plant and its controller at each level is a
\emph{lens}:
\begin{equation}
  \ell_i = (\text{get}_i, \text{put}_i)
  = (\obs_i : \X_i \to \Y_i, \quad
     \text{apply}_i : \U_i \times \X_i \to \X_i).
\end{equation}
The lens laws (get-put, put-get) formalize the requirement that observation
and control are consistent:
\begin{align}
  \text{put}(\text{get}(x), x) &= x
  & &\text{(get-put: no-op control preserves state)} \\
  \text{get}(\text{put}(u, x)) &= \obs(\dyn(x, u))
  & &\text{(put-get: effect is observable)}
\end{align}
\end{definition}


% =============================================================================
\section{EGRI as Level~2 Meta-Controller}
\label{sec:egri}
% =============================================================================

\begin{definition}[EGRI 9-tuple as $\rcs_2$]
\label{def:egri}
The Evaluator-Governed Recursive Improvement (EGRI) framework instantiates
$\rcs_2$ with:
\begin{equation}
  \rcs_2 = (\X_2, \Y_2, \U_2, \dyn_2, \obs_2, \shield_2, \ctrl_2)
\end{equation}
where:
\begin{center}
\begin{tabular}{@{}cll@{}}
  \toprule
  \textbf{RCS} & \textbf{EGRI} & \textbf{Meaning} \\
  \midrule
  $\X_2$ & Artifact state space $X$ & Controller parameters, architecture \\
  $\Y_2$ & Outcome $(\score, C)$ & Score vector + constraint satisfaction \\
  $\U_2$ & Mutation operators $\M$ & Proposed changes to $\X_2$ \\
  $\dyn_2$ & Trial execution & propose $\to$ execute $\to$ evaluate $\to$ select \\
  $\obs_2$ & Evaluator $\score$ & $\score : \text{traces} \to \setR^d$ \\
  $\shield_2$ & Immutable evaluator $+$ budget $+$ rollback & DP2, Law~2, DP3 \\
  $\ctrl_2$ & Proposer $+$ selector & Mutation strategy \\
  \bottomrule
\end{tabular}
\end{center}
\end{definition}

\paragraph{Two laws admit direct formal content.}
Among the Five Laws of EGRI, only two have derivations inside the
stability budget of Definition~\ref{def:budget}:

\begin{law}[Mutation--evaluation proportionality]
\label{law:capacity}
Do not grant an agent more mutation freedom than its evaluator can
reliably judge.  Formally, $\abs{\M} \leq \text{capacity}(\score)$
(Ashby's Law of Requisite
Variety~\cite{ashby1956introduction}).
\end{law}

\begin{law}[Budget closure]
\label{law:budget}
The loop must fail closed when budget is exhausted.  Formally, the
budget $\budget(t)$ is a Lyapunov function on the meta-controller:
\begin{equation}
  \budget(t_{k+1}) < \budget(t_k)
  \quad \text{and} \quad
  \budget(t_k) = 0 \implies \text{halt}.
\end{equation}
\end{law}

\paragraph{Three design principles.}
The remaining three ``Laws'' of EGRI are \emph{design principles}: they
are policy statements that shape the meta-controller but do not follow
from the stability budget and are not invoked in any proof in this
paper.

\textbf{Design Principle~1 (Evaluator supremacy).}\quad An optimization
loop is only as safe as the evaluator that governs it.

\textbf{Design Principle~2 (Evaluator immutability).}\quad Never mutate
the evaluator and the artifact in the same trial.  The analogy with the
separation principle of linear-quadratic-Gaussian control
(Wonham 1968) is motivating; EGRI operates on arbitrary artifacts, so
no LQG-level theorem follows, and this is kept as a principle rather
than a derived condition.

\textbf{Design Principle~3 (Rollback guarantee).}\quad Every promoted
state must be recoverable.  The set of promoted states is required, by
design, to be forward-invariant under the rollback operation; this
defines the rollback operator rather than being derived from earlier
assumptions.

\begin{proposition}[EGRI stability coupling]
\label{prop:egri-coupling}
When EGRI ($\rcs_2$) mutates the parameters of the autonomic controller
($\rcs_1$), the perturbation enters the Level~1 stability budget
via the design evolution term $L_{d,1} \, \eta_1$, where
$\eta_1 = \sup_k \norm{\Delta \param_{1,k}}$ is the maximum per-trial
parameter change.

For $\stab_1 > 0$ to hold, the EGRI mutation magnitude must satisfy:
\begin{equation}
  \eta_1 < \frac{\decay_1 - L_{\param,1}\rho_1 - \beta_1\taubar_1
  - \frac{\ln\nu_1}{\taua[1]}}{L_{d,1}}.
\end{equation}
\end{proposition}


% =============================================================================
\section{Observer: Fold as Sufficient Statistic}
\label{sec:observer}
% =============================================================================

\begin{definition}[Fold observer]
\label{def:fold}
At level $\Li{1}$, the \emph{fold observer} is a deterministic Mealy
accumulator over the event stream:
\begin{equation}
  \homostate_{k+1} = \text{fold}(\homostate_k, y_k),
\end{equation}
where $\homostate \in \X_1$ is the homeostatic state and $y_k \in \Y_1$
is the $k$-th event.  The fold maintains~$\homostate$ without
autonomous dynamics between events.  The analogy with a Luenberger
observer (with $A = I$ and full gain on measurements) holds only when
plant dynamics are slower than the folding rate~--- a regime in which
the identity transition absorbs all plant drift between successive
events.  Outside that regime, the Luenberger framing is notational, not
technical: no plant-prediction cancellation is performed.
\end{definition}

\begin{proposition}[Sufficiency]
\label{prop:sufficient}
If the controller $\ctrl_1$ at level $\Li{1}$ depends on the event history
$y_0, \ldots, y_k$ only through $\homostate_k$:
\begin{equation}
  \ctrl_1(y_0, \ldots, y_k) = \ctrl_1(\homostate_k),
\end{equation}
then $\homostate_k$ is a \emph{sufficient statistic} for regulation decisions
at level $\Li{1}$ with respect to the observation history.
\end{proposition}

\begin{remark}[Self-observation cost]
The agent measuring its own state consumes resources (tokens, context window
space), introducing a reflexive delay $\taubar_{\text{self}}$ that enters the
stability budget at level $\Li{1}$:
\begin{equation}
  \stab_1 = \decay_1 - L_{\param,1}\rho_1 - L_{d,1}\eta_1
  - \beta_1(\taubar_1 + \taubar_{\text{self}})
  - \frac{\ln\nu_1}{\taua[1]}.
\end{equation}
The hysteresis gate (deadband between activation and deactivation thresholds
with minimum dwell time) prevents self-referential oscillation by ensuring
small perturbations from introspection do not trigger state transitions.
\end{remark}


% =============================================================================
\section{Multi-Agent Extension}
\label{sec:fleet}
% =============================================================================

\begin{definition}[Fleet RCS]
\label{def:fleet}
A \emph{fleet} of $N$ agents is a composition of RCS instances sharing a
communication graph $G = (V, E)$:
\begin{equation}
  \rcs_{\text{fleet}} = \bigotimes_{G} \rcs^{(j)},
  \quad j = 1, \ldots, N,
\end{equation}
where $\otimes_G$ denotes composition with information exchange along edges
of $G$.  Each $\rcs^{(j)}$ is a full RCS hierarchy ($\Li{0}$ through
$\Li{3}$); the fleet adds a coordination layer governed by population
dynamics~\cite{quijano2017population}.
\end{definition}

\begin{definition}[Cooperative resilience~{\cite{chacon2025cooperative}}]
\label{def:resilience}
The \emph{cooperative resilience score} for variable $j$ under disruption $l$
is:
\begin{equation}
  J_{jl} = \frac{t_i + F_{jl} \cdot \Delta t_f + G_{jl} \cdot \Delta t_r}
  {t_i + \Delta t_f + \Delta t_r},
\end{equation}
where $F_{jl}$ and $G_{jl}$ are the failure and recovery profile integrals
respectively.  Fleet resilience is the \emph{harmonic mean} across all
variables, penalizing worst-case degradation:
\begin{equation}
  J_{\text{fleet}} = \left(\frac{1}{K} \sum_{j=1}^{K} J_j^{-1}\right)^{-1}.
\end{equation}
\end{definition}


% =============================================================================
\section{Self-Referential Closure}
\label{sec:closure}
% =============================================================================

The RCS hierarchy possesses a fixed-point property: Level~3 (governance) includes
the formalization itself.  The CLAUDE.md invariants, AGENTS.md operational rules,
METALAYER.md control manifest, and this very document are not external to the
system --- they \emph{are} the Level~3 controller $\ctrl_3$.  This section makes
the self-referential structure precise and verifies that it does not break
stability.

\begin{definition}[Self-referential RCS]
\label{def:self-ref}
An RCS hierarchy $\rcs_0, \ldots, \rcs_N$ has \emph{self-referential closure}
if the Level~$N$ controller $\ctrl_N$ includes a model of the complete hierarchy
$\{\rcs_0, \ldots, \rcs_N\}$.  That is, the governance rules, documentation,
and formalization at $\Li{N}$ \emph{are} the controller $\ctrl_N$:
\begin{equation}
  \ctrl_N \supseteq \text{Model}(\rcs_0, \ldots, \rcs_N).
\end{equation}
\end{definition}

\begin{proposition}[Governance stability]
\label{prop:governance}
Under the standard assumptions (Assumptions~\ref{ass:per-level}--\ref{ass:jump}),
Level~3 of the Life Agent OS satisfies $\stab_3 > 0$ with the following
concrete parameter values, derived from the Life runtime defaults and
cached in \texttt{data/parameters.toml} (the canonical source for all
numerics in this paper):
\begin{center}
\begin{tabular}{@{}cp{6.8cm}c@{}}
  \toprule
  \textbf{Parameter} & \textbf{Interpretation} & \textbf{Value} \\
  \midrule
  $\decay_3$ & Nominal decay --- policy crystallization
    (${\sim}\,1/100$ days) & $\rcsgammaLthreeDisp$ \\
  $L_{\param,3} \rho_3$ & Adaptation cost --- setpoint adjustment
    effectively zero at day scale & $\rcsadaptcostLthreeDisp$ \\
  $L_{d,3} \eta_3$ & Design cost --- structural changes
    (${\lesssim}\,1/\text{quarter}$) & $\rcsdesigncostLthreeDisp$ \\
  $\beta_3 \taubar_3$ & Delay cost --- human review latency
    (${\sim}\,1\,\text{hour}$) \textbf{dominates} & $\rcsdelaycostLthreeDisp$ \\
  $\frac{\ln \nu_3}{\taua[3]}$ & Switching cost --- profile switches
    once per day & $\rcsswitchcostLthreeDisp$ \\
  \midrule
  $\stab_3$ & Stability budget at Level~3
    & $\rcsmarginLthreeDisp > 0\;\checkmark$ \\
  \bottomrule
\end{tabular}
\end{center}
The governance margin is \emph{delay-dominated}: $\beta_3 \taubar_3$
accounts for ${\sim}56\%$ of the total cost, while the other three
perturbations are four orders of magnitude smaller.  This is the expected
regime at the top of the hierarchy --- governance operates on a day
scale, so adaptation, design evolution, and switching are effectively
static.  Stability is determined almost entirely by whether policy
crystallization ($\decay_3$) outpaces human review latency.
\end{proposition}

\begin{remark}[Narrowest margin]
\label{rem:narrowest}
Level~3 has the narrowest stability margin in the hierarchy
($\stab_3 = \rcsmarginLthreeDisp$ vs.\ $\stab_0 = \rcsmarginLzeroDisp$,
$\stab_1 = \rcsmarginLoneDisp$, $\stab_2 = \rcsmarginLtwoDisp$).
This is expected from the time-scale separation: governance operates on
the slowest time scale (days--weeks), so $\decay_3$ is smallest.  The
system is most fragile at its highest level --- rapid governance changes
would violate stability.  This matches the design intent: policy changes
should be rare and deliberate.
\end{remark}

\begin{remark}[Executable witnesses: numerical cross-check across representations]
\label{rem:witnesses}
The Rust traits \texttt{RecursiveControlledSystem\textless L\textgreater},
\texttt{LyapunovCandidate\textless L\textgreater}, and
\texttt{StabilityBudget} in the \texttt{aios-protocol} crate implement
the type signatures of Definitions~\ref{def:cs}--\ref{def:budget}, and
the Python tests verify the concrete parameter values of
Proposition~\ref{prop:governance}.  The agreement between this
document, the Rust types, and the numerical tests is best described as
a \emph{numerical cross-check across representations}: the same scalar
$\stab_3 = \rcsmarginLthreeDisp$ is recovered in three independent
encodings.  A semantic-preservation equivalence across the three
layers~--- in the sense of compiler verification~--- is not claimed.
\end{remark}


% =============================================================================
\section{Instantiation Catalogue}
\label{sec:instantiations}
% =============================================================================

The RCS 7-tuple is a \emph{type signature}.  Each concrete controller
populates the signature by discharging the per-level assumptions
(H1)--(H5) of Section~\ref{sec:stability} on its own state space.  We
illustrate the breadth of the signature with four instantiations, each
drawn from a distinct control-theoretic or learning-theoretic
community.  For every row we name $(\X, \dyn, \shield)$ concretely and
cite the result that produces the per-level~$\decay_i > 0$ required by
(H2).  The remaining standing hypotheses either follow directly
(sandwich bounds from the state-space structure, jump comparability
from hybrid-system conventions) or are discharged by the construction
of the next level above.

\begin{table}[h]
\centering
\small
\setlength{\tabcolsep}{4pt}
\begin{tabular}{@{}p{2.4cm}p{2.6cm}p{2.6cm}p{2.6cm}p{4.4cm}@{}}
\toprule
\textbf{Instantiation} & \textbf{State $\X$} & \textbf{Dynamics $\dyn$}
  & \textbf{Shield $\shield$} & \textbf{Decay witness for $\decay_i > 0$}\\
\midrule
LQR / classical
  & $\setR^n$ continuous
  & $Ax + Bu$ or smooth $\dyn(x,u)$
  & CBF / control-invariant set~\cite{ames2019cbf}
  & Riccati or LMI feasibility; quadratic $V = x^\top Px$~\cite{khalil2002nonlinear}.\\
\addlinespace
Dreamer actor--critic
  & Latent $z \in \setR^d$ from RSSM
  & Trained Lipschitz RSSM predictor
  & Barrier on imagined trajectories
  & Spectral-norm bound of the trained RSSM under symlog + KL-balance yields bounded latent drift~\cite{hafner2025dreamerv3}.\\
\addlinespace
V-JEPA~2 latent
  & JEPA latent $\hat z \in \setR^d$
  & JEPA predictor + CEM-MPC
  & Intrinsic energy cost
  & Bounded prediction error over horizon $T \leq 16$ frames gives finite-horizon Lyapunov decay of the goal-conditioned energy~\cite{assran2025vjepa2}.\\
\addlinespace
LLM-as-controller
  & Context window + meaning-space lift
  & Bhargava stochastic kernel
  & AgentSpec PDA~\cite{wang2026agentspec}
  & Attention-block $\sigma_{\max}$~\cite{bhargava2023magic} plus meaning-space quotient of~\cite{soatto2023taming} yields reachability + lifted contraction sufficient for (H2).\\
\bottomrule
\end{tabular}
\caption{Four instantiations of the RCS type signature.  Each row
populates the 7-tuple on its own state space; the decay-witness column
produces (H2).  The remaining per-level hypotheses (H1, H3--H5) follow
from each row's state-space conventions (sandwich bounds,
adaptation/design Lipschitz bounds, dwell-time structure); the global
conditions (H6)--(H7) are imposed on the composed hierarchy.  When (H1,
H3--H7) also hold, Theorem~\ref{thm:recursive} composes the rows into a
finite-depth hierarchy with $\omega_N > 0$.}
\label{tab:instantiations}
\end{table}

Two observations on Table~\ref{tab:instantiations}.  First, the
witnesses cited for the LLM row establish reachability and lifted
contraction, not a quadratic Lyapunov bound in the raw token space; the
per-level assumption is therefore discharged in the meaning-space
quotient of~\cite{soatto2023taming}, not in the raw context window.
Second, the table is not exhaustive: additional instantiations
(PID-style activation steering, CBF-QP over continuous plants, EGRI as
$\rcs_2$) can be added by identifying their state space, dynamics,
shield, and per-level decay witness.


% =============================================================================
% BIBLIOGRAPHY
% =============================================================================
\bibliographystyle{plainnat}
\bibliography{\SHAREDLATEX/references}

\end{document}
